\documentclass{situationist-preprint}

\title{The Man in the Room Is Not the Subject}
\msid{MND-07-b}
\msfamily{MND-07}
\msdiscipline{Philosophy of Mind}
\mstype{Companion manuscript}
\mskeywords{Chinese Room; systems reply; subject of understanding; executing; running}

\begin{document}
\maketitle

\begin{abstract}
Searle's Chinese Room asks us to consult the man in the room about whether he understands Chinese. We argue that the man is not the subject to whom understanding of Chinese should be ascribed, and that his testimony is therefore not evidence about it. The subject of a cognitive property is whatever has the functional organization in virtue of which the property is had. In the room, that is the system that runs the program, not the man who executes its steps. The man's sincere denial is true of him and silent about the system.
\end{abstract}

\section{Which subject?}

The Chinese Room has two candidates for the subject of understanding: the man, who carries out the rules, and the system comprising the man, the rule book and the slips of paper \citep{Searle1980}. MND-07-a distinguishes executing a program, carrying out its operations one by one, from running it, having states and transitions that realize the functional organization the program specifies (MND-07-a, Definitions~2.1 and 2.2). The man executes the program; the system runs it.

\begin{principle}[Subject of a functional property]\label{p:subject}
If a cognitive property is had in virtue of a functional organization, its subject is the bearer of that organization: the entity whose states play the roles the organization specifies.
\end{principle}

\begin{proposition}\label{p:notman}
If understanding Chinese is had in virtue of a functional organization, the man in the room is not its subject.
\end{proposition}

\begin{proof}
The man's states are beliefs and intentions about which characters to write after which; they do not play the roles of beliefs about what the Chinese sentences are about. The states that play those roles are configurations of the man together with the book and the slips, which belong to the system. By Principle~\ref{p:subject}, the subject of the property, if anything is, is the system.
\end{proof}

The antecedent is the functionalism of Assumption~4.1 of MND-07-a. Without it, Proposition~\ref{p:notman} still has a weaker form: Searle's argument gives no reason to think the man is the relevant subject, since the inference from the man's lack of understanding to the system's requires the downward transfer principle that MND-07-a shows false for other cognitive properties (Principle~P1 and Proposition~3.1).

\section{The man's testimony}

Searle's case draws its force from first-person testimony: the man sincerely reports that he understands nothing. The testimony is reliable, and it is true of him. But testimony concerns the subject who gives it. The man's report that he does not understand Chinese is evidence about the man, as a neuron's failure to respond to a meaning, if neurons could report, would be evidence about the neuron. Asked in Chinese, the system answers about the topics of the conversation; asked in English, the man denies knowing anything about them. These are reports of two subjects with different states, not conflicting reports of one.

\begin{objection}
The system is nothing over and above the man and some paper. Paper does not understand, so if the man does not, nothing in the room does.
\end{objection}

\begin{reply}
The objection relies on downward transfer in reverse: that a system has a cognitive property only if one of its components has it. Brains are counterexamples. The system is not an extra object alongside the man and the paper, but it has a functional organization that none of them has, and Principle~\ref{p:subject} locates the property there. When the man internalizes the rule book, the same holds: he executes the program from memory, and a system realized in his brain, but not identical to him or to any part of it, runs it (MND-07-c, Proposition~1.1; compare \citealp{Cole1991}).
\end{reply}

\section{Conclusion}

The subject of a functional cognitive property is the bearer of the organization in virtue of which it is had. In the Chinese Room, that is the system. The man executes the program and does not understand Chinese; he is not the subject of which understanding is in question, and his testimony, true of him, says nothing about it.

\begin{thebibliography}{9}\small
\bibitem[Searle(1980)]{Searle1980} Searle, J.~R. (1980). Minds, brains, and programs. \emph{Behavioral and Brain Sciences}, 3(3), 417--457.
\bibitem[Cole(1991)]{Cole1991} Cole, D. (1991). Artificial intelligence and personal identity. \emph{Synthese}, 88(3), 399--417.
\end{thebibliography}

\end{document}
